\SourceLeaf{255}%
from which we readily conclude that equations (5) and (6) are entirely sufficient for our case, provided that what they leave undetermined is chosen so that \(r\) and \(\varphi\) agree with the required values at the initial point and at points at an infinitesimal distance from it.

As for the integration of equations (5) and (6) itself, it is known that it can be reduced to the integration of ordinary differential equations. These generally become so intricate that little is gained. By contrast, expansion in series presents no difficulty, and the series are more than sufficient for practical use whenever the portions of the surface in question are small. The formulas given thus offer a rich source for solving many problems of great importance. Here, however, we shall work out just one example to show the nature of the method.
\Article{23}
We shall consider the case in which all lines along which \(p\) is constant are shortest lines that intersect the line given by \(\varphi=0\) at right angles. We can regard the latter as the line of abscissas. Let \(A\) be the point for which \(r=0\), let \(D\) be any point on the line of abscissas, and let \(AD=p\). Let \(B\) be any point on the shortest line normal to \(AD\) at \(D\), and let \(BD=q\). Then \(p\) can be regarded as the abscissa and \(q\) as the ordinate of \(B\). We take the abscissas as positive on the branch of the line of abscissas corresponding to \(\varphi=0\), while we always regard \(r\) as positive. We take the ordinates as positive on the side where \(\varphi\) is measured between 0 and \(180^0\).

By the theorem of Article 16 we shall have \(\omega=90^0\), \(F=0\). We also have \(G=1\), and we set \(\BareRoot{E}=n\). Thus \(n\) will be a function of \(p,q\). For \(q=0\) it must be \(=1\). Applying the formula given in Article 18 to our case shows that on \textit{any} shortest line we must have \(\mathrm{d}\theta=-\frac{\mathrm{d}n}{\mathrm{d}q}\mathord{.}\mathrm{d}p\), where \(\theta\) denotes the angle between an element of this line and the element of the line along which \(q\) is constant. Now the line of abscissas is itself a shortest line, and along it \(\theta=0\) everywhere. Therefore, for \(q=0\), we must have \(\frac{\mathrm{d}n}{\mathrm{d}q}=0\) everywhere. Hence if we expand \(n\) in a series in ascending powers of \(q\), the series must have the following form
\[
n=1+fqq+gq^3+hq^4+\text{etc.}
\]
where \(f,g,h\) and so on are functions of \(p\). We set
